\cotefile{71}
\scanpage{38}\mspage{$\underline{\overline{\mathrm{XI}}}$}
3 - \uline{Preuve du \unsure{théorĕme|théorème}} 1

\margin{!}Soit $\Gamma$ un polygone et $E(\Gamma)$ son enveloppe convexe. Soit $s,s'\in S(\Gamma)$ un couple de sommets \deleted{\illegible{1}} tels que $s'$ est $(\mathbb{R}^2-\Gamma)$-visible de $s$ et $]s,s'[\subset\partial E(\Gamma)$ (cf. coroll. au lemme $\underline{\overline{\mathrm{VIII}}}$). Soient $\Gamma^1_{ss'}$ et $\Gamma^2_{ss'}$ les deux arcs de polygone \inserted{(de $\Gamma$)} réliant $s$ à $s'$ (cf. Remarque 1).
\margin{$\vert$}

Soit $\Lambda_{ss'}$ l'arc de polygone $[s,s']$. On appelle $\Lambda^1$ (resp. $\Lambda^2$) le polygone obtenu en récollant $\Lambda_{ss'}$ avec $\Gamma^1_{ss'}$ (resp. $\Gamma^2_{ss'}$) selon $\{s,s'\}$.

% slanted
\margin{\shortstack{$a=s$\\$b=s'$}\,?}
\uline{Lemme} $\underline{\overline{\mathrm{XII}}}$. Si $\Lambda^1$ (resp. $\Lambda^2$) est un $J$-polygone alors $\Gamma^2_{ab}-\{a,b\}\subset\operatorname{int}\Lambda^1$ ou bien $\Gamma^2_{ab}-\{a,b\}\subset\operatorname{ext}\Lambda^{\deleted{\unsure{2}}1}$ (resp. $\Gamma^1_{ab}-\{a,b\}\subset\operatorname{int}\Lambda^2$ $\Gamma^1_{ab}-\{a,b\}\subset\operatorname{ext}\Lambda^2$)
\margin{1}

\uline{Demonstration} : Il est clair que $\Gamma^2_{ab}-\{a,b\}\subset\deleted{\unsure{\partial|E}}\mathbb{R}^2-\Lambda^1$ car $]s,s'[\cap\Gamma=\varnothing$, mais en plus $\Gamma^2_{ab}-\{a,b\}$ est connexe donc la conclusion. On réfait l'argument pour $\Lambda^2$, c.q.f.d.
\margin{$\vert$}\margin{$\vert$}

On suppose maintenant \inserted{que} $\Lambda^1$ et $\Lambda^2$ \deleted{\unsure{que}} sont des $J$-polygones.

\uline{Lemme} $\underline{\overline{\mathrm{XIII}}}$ : Si $x\in\partial E(\Gamma)$ et $D_x$ est un disque de centre $x$ \unsure{vérifiont|vérifiant} b) de la definition 2 par rapport à $\partial E(\Gamma)$ on a alors :
\[
\operatorname{ext}\partial E(\Gamma)\cap\deleted{D_x}\subset\operatorname{ext}\Lambda^i,\qquad i=1,2.
\]
\margin{\tikz[baseline=(ex.base)]{\node[inner sep=3pt] (ex) {$\begin{aligned}
\operatorname{Ext}(\partial E(\Gamma))&=\mathbb{R}^{\unsure{2|7}}-E(\Gamma)\\
&\subset\operatorname{Ext}(\Lambda^i)
\end{aligned}$};
\draw (ex.north east) to[bend right=3] (ex.north west) -- (ex.south west) to[bend left=3] (ex.south);}}

\uline{Démonstration} : Comme $E(\Lambda^i)\subset E(\Gamma)$, $i=1,2$, on a $\operatorname{ext}\deleted{\partial E(}\Lambda^i\deleted{)}\supset\operatorname{ext}\partial E(\Lambda^i)\supset\operatorname{ext}\partial E(\Lambda)$ ce qui entraîne l'ennoncé.
% archive sheet number: 37

\scanpage{39}\mspage{\deleted{\unsure{X bis|XI}}$\underline{\overline{\mathrm{XII}}}$}
On observe que $\Lambda^1$ et $\Lambda^2$ ont une arète en commun\inserted{, savoir,} $[s,s']$ et $\Lambda^1\cap\Lambda^2=[s,s']$.

\uline{Lemme} XIV. Si $D_s$ est un disque satisfaisant b) de la definition 2 par rapport à $\Lambda^1$ et à $\Lambda^2$ on a que $\operatorname{int}\Lambda^1\supset\Lambda^2\cap D_s-[s,s']\cap D_s$ ou bien $\operatorname{int}\Lambda^2\supset\Lambda^1\cap D_s\deleted{\unsure{-|\subset}}[s,s']\cap D_s$ \deleted{\unsure{(}}et resp. \inserted{que} $\operatorname{ext}\Lambda^2\supset\Lambda^1\cap D_s-[s,s']\cap D_s$ ou bien $\operatorname{ext}\Lambda^1\supset\Lambda^2\cap D_s-[s,s']\cap D_s$\deleted{)}.

\uline{Demonstration} : D'aprés le lemme XIII, $\operatorname{ext}\Lambda^1\cap D_s\supset\operatorname{ext}\partial E(\Gamma)\cap D_s$, $\operatorname{ext}\Lambda^2\cap D_s\supset\operatorname{ext}\partial E(\Gamma)\cap D_s$ donc $\operatorname{ext}\Lambda^1\cap\operatorname{ext}\Lambda^2\cap D_s\ne\varnothing$.

Comme $\Lambda^1\cap D_s$ et $\Lambda^2\cap D_s$ ont un coté en commun il résulte que le \inserted{secteur angulaire} $\operatorname{ext}\Lambda^1\cap D_s$ contient ou est contenu dans l'un des secteurs angulaires determinés par $\Lambda^2\cap D_s$. Mais $\operatorname{ext}\Lambda^1\cap\operatorname{ext}\Lambda^2\cap D_s\ne\varnothing$ entraîne que ce secteur est précisement $\operatorname{ext}\Lambda^2\cap D_s$, cela implique (sans peine)
\[
\operatorname{int}\Lambda^1\supset\Lambda^2\cap D_s-[s,s']\cap D_s
\quad\text{ou bien}\quad
\deleted{\operatorname{ext}}\operatorname{int}\Lambda^2\supset\Lambda^1\cap D_s-[s,s']\cap D_s,
\]
c.q.f.d.

La conjonction du lemme $\underline{\overline{\mathrm{XII}}}$ et du lemme $\underline{\overline{\mathrm{XIII}}}$ \deleted{\unsure{entraîne|entraînent}} montre que
\[
\begin{aligned}
\operatorname{int}\Lambda^1\supset{}&\Lambda^2-\{g,s'\}
&&\text{ou bien }\operatorname{int}\Lambda^2\supset\Lambda^1-\{s,s'\}\ \text{et on a}\\[-2pt]
&\inserted{\substack{\text{''}\\\Gamma^2_{s,s'}-\{s,s'\}}}
&&\inserted{-\qquad\substack{\text{''}\\\Gamma^{\unsure{1|2}}_{s,s'}.-\{s,s'\}}}\\
\text{respectivement}\quad&\operatorname{ext}\Lambda^2\supset\Lambda^1-\{s,s'\}
&&\text{et }\operatorname{ext}\Lambda^1\supset\Lambda^2-\{s,s'\},
\end{aligned}
\]

Supposons pour fixer les idées que $\operatorname{int}\Lambda^1\supset\Lambda^2-\{s,s'\}$ donc aussi que $\operatorname{ext}\Lambda^2\supset\Lambda^1-\{s,s'\}$.

On a donc :

\uline{Lemme XV} : $\operatorname{int}\Lambda^2\subset\operatorname{int}\Lambda^1$.
% archive sheet number: 38

\scanpage{40}\mspage{$\underline{\overline{\mathrm{XIII}}}$}
\uline{Demonstration} : Pour voir que $\operatorname{int}\Lambda^2\subset\deleted{\unsure{\operatorname{ext}}}\operatorname{int}\Lambda^1$ \inserted{on} observe d'abord que :
\[
\Lambda^1-\{s,s'\}\subset\operatorname{ext}\Lambda^2\quad\text{donc}\quad\deleted{\unsure{\operatorname{ext}}}\operatorname{int}\Lambda^2\cap(\Lambda^1-\{s,s'\})=\varnothing
\]
, et par conséquent :
\[
\operatorname{int}\Lambda^2\subset\operatorname{int}\Lambda^1\quad\text{ou bien}\quad\operatorname{int}\Lambda^2\subset\operatorname{ext}\Lambda^1\quad(\text{car }\operatorname{int}\Lambda^2\subset\mathbb R^2-\Lambda^1).
\]
Mais $\operatorname{int.}\Lambda^2\cap\operatorname{int}\Lambda^1\ne\varnothing$. En effet soit $x\in{]s,s'[}$ et $D_x$ un disque de centre $x$ qui vérifie b) de la def.1 par rapport à $\Lambda^1$, $\Lambda^2$ et $\partial E(\Gamma)$, le lemme $\underline{\overline{\mathrm{XIII}}}$ montre que $D_x\cap\operatorname{ext}\partial E(\Gamma)=D_x\cap\operatorname{ext}\Lambda^1=D_x\cap\operatorname{ext}\Lambda^1$ donc $D_x\cap\deleted{\unsure{\operatorname{ext}}}\operatorname{int}\Lambda^1=D_x\cap\deleted{\unsure{\operatorname{ext}}}\operatorname{int}\Lambda^2$, et par conséquent $\operatorname{int}\Lambda^1\cap\operatorname{int}\Lambda^2\ne\varnothing$, c.q.f.d.

On veut prouver que $\Gamma$ est un $J$-polygone par récurrence sur $\operatorname{card}S(\Gamma)-2$.

Si $\operatorname{card}S(\Gamma)-2=1\quad\Rightarrow\quad\Gamma={}$ au bord d'un triangle\inserted{et le \unsure{theorēme|theorème} 1 devient donc trivial}, donc il suffit de prouver que si le \unsure{théorĕme|théorème} 1 est vrai pour tout $\Gamma$ avec $\operatorname{card}S(\Gamma)-2<n$ alors il est aussi vrai pour tout $\Gamma$ avec $\operatorname{card}S(\Gamma)-2=n$.

Il est clair que $\operatorname{card}S(\Lambda^i)<\operatorname{card}S(\Gamma)$ \inserted{$i=1,2$}, autrement $\Lambda^i=\Gamma$ ce qui est absurd puisque $ {]s,s'[}\cap\Gamma=\varnothing$. Les lemmes $\underline{\overline{\mathrm{XII}}}$ et $\underline{\overline{\mathrm{XIV}}}$ \inserted{avec le lemme $\underline{\overline{\mathrm{XV}}}$} nous permettent de supposer sans perte de generalité que : a) $\operatorname{int}\Lambda^2\subset\operatorname{int}\Lambda^1$ ;

b) $\Lambda^2-\{s,s'\}=\Gamma^2_{s,s'}-\{s,s'\}\subset\operatorname{int}\Lambda^1$.

Posons donc :
\[
\begin{aligned}
V&=\operatorname{ext}\Lambda^1\cup\operatorname{int}\Lambda^2\cup{]s,s'[}\quad\text{et}\\
U&=\operatorname{int}\Lambda^1-(\operatorname{int}\Lambda^2\cup(\Gamma^2_{ss'}\cap\operatorname{int}\Lambda^2)).
\end{aligned}
\]
Par \unsure{hypothēse|hypothèse} inductive $\operatorname{ext}\Lambda^1$ est non-relativement-compact donc $V$ non plus et $U\subset\operatorname{int}\Lambda^1$ est relativement compact puisque $\operatorname{int}\Lambda^1$ est relativement compact par hypothēse de récurrence.

\figure{\makebox[\linewidth][r]{\begin{tikzpicture}[x=1cm,y=1cm]
\fill (0,0) -- (.11,.015) -- (.15,-.06) -- cycle;
\draw (.05,-.015) .. controls (.94,-.34) and (1.60,-.80) .. (3.8,-1.20);
\end{tikzpicture}}}
% archive sheet number: 39

\scanpage{41}
% Two lines written upside down relative to the archive sheet number.
Demonstration : pour voir que $\operatorname{int}\Lambda^2\supset\operatorname{int}\Lambda^1$ il suffit de prouver que tout point de
% archive sheet number: 40

\scanpage{42}\mspage{$\underline{\overline{\mathrm{XIV}}}$}
\uline{Lemme} $\underline{\overline{\mathrm{XVI}}}$. On a : a) $\mathbb R^2-\Gamma=U\cup V$

b) $U$ et $V$ sont des ouverts connexes, $U$ relativement compact, $V$ non-relativement compact.

\uline{Démonstration} : a) On a $U\cup V=\operatorname{ext}\Lambda^1\cup\operatorname{int}\Lambda^1\cup{]s,s'[}-(\Gamma^2_{ss'}-\{s,s'\})$.

D'autre côté $\Gamma=\Lambda^1-{]s,s'[}\cup(\Gamma^2_{ss'}-\{s,s'\})$, mais alors
\[
\begin{aligned}
\mathbb R^2-\Gamma&=(\mathbb R^2-\Lambda^1)\cup{]s,s'[}-(\Gamma^2_{ss'}-\{s,s'\})=\operatorname{ext}\Lambda^1\cup\operatorname{int}\Lambda^1\cup{]s,s'[}\\
&\hspace{12em}-(\Gamma^2_{ss'}-\{s,s'\}).
\end{aligned}
\]

b) Prouvons que $U$ est connexe. Soient $x,y\in U$\unsure{;|,} montrons que $x$ peut être relié\deleted{\unsure{r|e}} à $y$ par un arc de polygone .contenu dans $U$. Si $x,y\in U$ alors $x,y\in\operatorname{int}\Lambda^1$.

On peut donc reli$\overset{\unsure{\prime|\cdot}}{\text{e}}$r $x$ à $y$ par un arc de polygone $\Lambda_{xy}\subset\operatorname{int}\Lambda^1$. Si $\Lambda_{xy}\subset\underbrace{\operatorname{int}\Lambda^1-R_{\Lambda^2}}_{\substack{\text{''}\\\underset{\cdot}{U}}}$ $\bigl(R_{\Lambda^2}=\operatorname{int}\Lambda^2\cup\Lambda^2\bigr)$ il n'y a plus rien à montrer.

Si $\Lambda_{xy}\cap(\Gamma^2_{ss'}-\{s,s'\})=\varnothing$ alors $\Lambda_{xy}\cap\operatorname{int}\Lambda^2=\varnothing$, puisque $x,y\in\operatorname{ext}\Lambda^2$ et $\Lambda_{xy}\cap\Lambda^2=\varnothing$ (car $\Lambda^2=[s,s']\cup(\Gamma^2_{ss'}-\{s,s'\})$ et $\Lambda_{xy}\cap[s,s']=\varnothing$ puisque $\Lambda_{xy}\subset\operatorname{int}\Lambda^{\unsure{4|1}}$)

Si $\Lambda_{\unsure{x}y}\cap(\Gamma^2_{ss'}-\{s,s'\})\ne\varnothing$ on peut trouver $x'$ et $y'\in(\Gamma^2_{ss'}-\{s,s'\})\cap\Lambda_{xy}$ tels que $\Lambda_{xx'}-\{x'\}\subset\operatorname{ext}\Lambda^2$, $\Lambda_{y'y}-\{y'\}\subset\operatorname{ext}\Lambda^2$ (ici $\Lambda_{xx'}$ (resp. $\Lambda_{y'y}$) dénote l'arc de polygone obtenu à partir de $\Lambda_{xy}$ qui rélie $x$ a $x'$ (resp. $y'$ à $y$) \inserted{obtenu} de foçon analogue a \unsure{$\overset{\unsure{2|\mathrm c|?}}{\text{ce}}$lui} de la Remarque 1), pour cela il suffit d'avoir
\[
\tikz[baseline=(eq.base)]{\node[inner sep=3pt] (eq) {$\begin{aligned}
\Lambda_{xx'}\cap(\Gamma^2_{ss'}-\{s,s'\})&=x',\\
\Lambda_{y y'}\cap(\Gamma^2_{ss'}-\{s,s'\})&=y'.
\end{aligned}$};
\draw ([xshift=2pt]eq.north west) -- (eq.south west) -- (eq.south east) -- ([xshift=4pt]eq.north east);
\node[inner sep=.2pt] at ([xshift=9pt]eq.east) {$\unsure{*}$};}
\]

\deleted{Soit maintenant $\unsure{\Gamma}$} \inserted{Comme} On a $x',y'\in\Gamma^2_{ss'}-\{s,s'\}$ on peut trouver un arc de polygone $\Gamma_{x'y'}$ de $\Gamma$, $\Gamma_{x'y'}\subset\Gamma^2_{ss'}-\{s,s'\}$ (cf. Remarque 1).
% archive sheet number: 41

\scanpage{43}\mspage{$\underline{\overline{\mathrm{XV}}}$}
Choisissons un voisinage $V$ de $\Gamma_{x'y'}$ tel que : $\mathrm{I})\ V\cap\Gamma_{ss'}=\varnothing$ ;
% Paired vertical strokes at the two edges of the upper sheet.
\reversemarginpar\margin{$\vert$}\normalmarginpar\margin{$\vert$}%
$\overline{\underline{\mathrm{II}}})\ V\cap\mathrm{int}\,\Lambda^{2}$, $V\cap\mathrm{ext}\,\Lambda^{2}$ sont des ouverts connexes. (Il \unsure{\.{y}|\'{y}|y} en a toujours car si $A$ et $B$ sont deux parties compactes de $\mathbb{R}^{2}$ il \unsure{exìste|existe} un couple de voisinages $A\subset V_A$, $B\subset V_B$ tel que $V_A\cap V_B=\varnothing$, le lemme I \deleted{\unsure{fai|fait}} montre qu'on peut trouver $V\inserted{\subset V_A}$ satisfaisant en plus $\overline{\underline{\mathrm{II}}})$ si $A=\Gamma_{\deleted{\unsure{s}}x'\deleted{\unsure{s}}y'}$ et $B=\Gamma_{ss'}$)

Comme $\Lambda_{xx'}\cap V\inserted{-\{x'\}}\ne\varnothing$ \deleted{\illegible{1}} (resp. $\Lambda_{y'y}\cap V-\{y'\}\ne\varnothing$) on a
\[
\begin{aligned}
&\Lambda_{xx'}\cap V-\underbrace{\Lambda_{xx'}\cap\Lambda^{2}}_{x'}\subset V\cap\mathrm{int}\,\Lambda^{2}\ \cup\ V\cap\mathrm{ext}\,\Lambda^{2}\quad(\text{resp.}\\
&\Lambda_{y'y}\cap V-\underbrace{\Lambda_{y'y}\cap\Lambda^{2}}_{y'}\subset V\cap\mathrm{int}\,\Lambda^{2}\ \cup\ V\cap\mathrm{ext}\,\Lambda^{2})\quad\text{mais alors}\\
&\Lambda_{xx'}\cap V-\{x'\}\subset V\cap\mathrm{ext}\,\Lambda^{2}\quad(\text{resp. }\Lambda_{y'y}\cap V-\{y'\}\subset V\cap\mathrm{ext}\,\Lambda^{2})\\
&\text{car}\quad \Lambda_{xx'}-\{x'\}\subset\mathrm{ext}\,\Lambda^{2}\quad(\text{resp. }\Lambda_{y'y}-\{y'\}\subset\mathrm{ext}\,\Lambda^{2}).
\end{aligned}
\]
Soit donc $x'',\deleted{y''\in}\ \in(\Lambda_{xx'}-\{x'\})\cap V$ (resp. $y''\in(\Lambda_{y'y}\unsure{\leftarrow|-}\{y'\})\cap V$).
Il est clair qu'on peut rélier $x''$ à $y''$ par un arc de polygone $\inserted{\Lambda_{x''y''}}\subset V$ avec $\Lambda_{x''y''}\cap\Gamma^{\prime}_{ss'}=\varnothing$ (cela résulte \unsure{dee|deu|du} fait que $V\cap\Gamma'_{ss'}=\varnothing$) alors l'arc de polygone $\inserted{\overline{\Lambda}_{xy}}$ obtenu
\reversemarginpar\margin{$\vert$}\normalmarginpar\margin{$\vert$}%
en récollant $\Lambda_{xx''}$ avec $\Lambda_{x''y''}$, $\Lambda_{x''y''}$ avec $\Lambda_{y''y}$, selon $\{x''\}$ et $\{y''\}$ respectivement, satisfait $\overline{\Lambda}_{x\dot y}\subset U=\unsure{\mathrm{Int}|\mathrm{int}}\,\Lambda^{1}-R_{\Lambda^{2}}\cap\mathrm{int}\,\Lambda^{1}$.

Prouvons que $V$ est connexe. Soient $x,y\in U=\mathrm{ext}\,\Lambda^{1}\cup\mathrm{int}\,\Lambda^{2}\cup{]s,s'[}$

Si $x,y\in\mathrm{ext}\,\Lambda^{1}$ on rélie $\dot{x}$ à $y$ par un arc de polygone $\Lambda_{xy}\subset\mathrm{ext}\,\Lambda^{1}$.

Si $x,y\in\mathrm{int}\,\Lambda^{2}\cup{]s,s'[}$ on peut aussi rélier $x$ à $y$ par un arc de polygone $\Lambda_{xy}\subset\mathrm{ext}\,\Lambda^{1}$.

Donc il reste le cas ou $x\in\mathrm{ext}\,\Lambda^{1}$ et $y\in\mathrm{int}\,\Lambda^{2}\cup{]s,s'[}$. On choisit $z\in{]s,s'[}$ on rélie $x$ a $z$ par $\Lambda_{xz}$ tel que $\Lambda_{xz}-\{z\}\subset\mathrm{ext}\,\Lambda^{1}$ et $y$ à $z$ par $\Lambda_{zy}$ tel que $\Lambda_{zy}\deleted{\unsure{-\{z\}}}\subset\mathrm{int}\,\Lambda^{2}\cup{]s,s'[}$. Il est \deleted{\unsure{q}}clair \unsure{quén|qu'en} récollant $\Lambda_{xz}$ avec $\Lambda_{zz}$ selon $\{z\}$ on obtient $\Lambda_{xy}\subset V$.

On a vu que $U$ et $V$ sont des connexes il sont necessairement ouverts car il sont les composantes connexes de $\mathbb{R}^{2}-\Gamma$, et $\Gamma$ est compact donc fermé dans $\mathbb{R}^{2}$\unsure{)|,|;}

\hfill c.q.s.d.
% archive sheet number: 42

\scanpage{44}\mspage{$\underline{\overline{\mathrm{XVI}}}$}
\uline{Lemme} $\underline{\overline{\mathrm{XVII}}}$. Si $x\in\Gamma$ et $D_x$ est un disque de centre $x$ qui ne coupe que les arêtes auxquelles $x$ appartient alors $D_x\cap U\ne\varnothing$, $D_x\cap V\ne\varnothing$.

\uline{Démonstration} : D'aprés la remarque qui suit la def. 1 il suffit d'exhiber un disque $D_x$ qui ne coupe que les arètes auxquelles $x$ appartient et tel que $D_x\cap U\ne\varnothing$, $D_x\cap V\ne\varnothing$

Si $x\in\Gamma^{1}_{ss'}-\{s,s'\}$ soit $D_x$ un disque de centre $x$ vérifiant : a) $D_x\subset\mathbb{R}^{2}-\Lambda^{2}$ ; b) $D_x$ satisfait b) de la \unsure{definition|defnition} 1 par rapport à $\Lambda^{1}$, donc $D_x\cap\mathrm{ext}\,\Lambda^{1}\ne\varnothing$ et par conséquent $D_x\cap V\ne\varnothing$. Aussi $D_x\subset\mathrm{ext}\,\Lambda^{2}$ et $D_x\cap\mathrm{int}\,\Lambda^{1}\ne\varnothing$ entraîne $D_x\cap V\ne\varnothing$. Si $x\in\Gamma^{2}_{ss'}-\{s,s'\}$ un argument analogue montre que $D_x\cap U\ne\varnothing$ et $D_x\cap V\ne\varnothing$.

Soit maintenant un disque $D_s$ qui satisfait b) de la \unsure{definition|defnition} 1 par rapport a $\Lambda^{1}$ et à $\Lambda^{2}$ on sait par le lemme \inserted{\textcircled{\scriptsize$\underline{\overline{\mathrm{XIV}}}$}} que
\[
\begin{aligned}
D_s\cap{}&\mathrm{ext}\,\Lambda^{1}\ne\varnothing
  &\quad\text{et}\quad D_s\cap{}&(\mathrm{int}\,\Lambda^{1}-\mathrm{int}\,\Lambda^{2})\\[-.5ex]
 &\hphantom{\mathrm{e}}\inserted{\begin{gathered}[t]\cap\\V\end{gathered}}
  &\begin{gathered}[t]\vphantom{\cap}\\\ne\varnothing\end{gathered}
  &\hphantom{(\mathrm{int}\,\Lambda^{1}-}\inserted{\begin{gathered}[t]\cap\\U\end{gathered}}
\end{aligned}
\]
ce qui suffit à prouver le lemme s$\dot{\mathrm{p}}$our $s$, un argument analogue démontre le lemme pour $s'$.

La \unsure{conjonetion|conjonction} du lemme $\underline{\overline{\mathrm{XVI}}}$ et du lemme $\underline{\overline{\mathrm{XVII}}}$ montre le thēorēme 1.

4- \uline{Preuve du théorēme 2} (On suppose que tout \inserted{polygone} $\Gamma$ est un $\unsure{J|j|y}$-polygone)

On procède par récurrence sur $\deleted{\unsure{S|s}}\mathop{\mathrm{card}} S(\Gamma)-2$.

On choisit $s,s'\in S(\Gamma)$ satisfaisant la condition du lemme XI, i.e. $s'$ est visible de $s$ de l'interieur.

\deleted{\unsure{Posons}} \inserted{Soient} $\Gamma^{1}_{ss'}$ et $\Gamma^{2}_{ss'}$ \deleted{\unsure{égale|égals}} les \unsure{deux|dex} arcs de polygone de $\Gamma$ réliant $s$ à $s'$ décrits dans la Remarque 1, et soit $\Lambda_{ss'}$ l'arc de polygone $\Lambda_{ss'}=[s,s']$. Appelons $\Lambda^{1}$ \deleted{l'arc} (resp. $\Lambda^{2}$) le polygone \deleted{f}obtenu à partir de $\Gamma^{1}_{ss'}$ (resp. $\Gamma^{2}_{ss'}$) et $\Lambda_{ss'}$ par \unsure{rēcollement|récollement} selon $\{s,s'\}$.

Il est clair que $\inserted{\mathrm{card}}\,S(\Lambda^{1})<\mathop{\mathrm{card}} S(\Gamma)$ (resp. $\mathop{\mathrm{card}} S(\Lambda^{2})<\mathop{\mathrm{card}} S(\Gamma)$), autre\textsuperscript{.}\\ment
% archive sheet number: 43

\scanpage{45}\mspage{$\underline{\overline{\mathrm{XVII}}}$}
$[s,s']\subset A(\Lambda^{1})$ (resp. $A(\Lambda^{2})$) ce qui est absurde \unsure{puìsque|puisque}
\[
{]s,s'[}\subset\mathbb{R}^{2}-\Lambda^{1}\quad(\text{resp. }\mathbb{R}^{2}-\Lambda^{2}).
\]

\uline{Lemme} $\underline{\overline{\mathrm{XVIII}}}$. $\Lambda^{2}-\{s,s'\}\deleted{\unsure{\supset}}\subset\mathrm{ext}\,\Lambda^{1}$ et $\Lambda^{1}-\{s,s'\}\subset\mathrm{ext}\,\Lambda^{2}$.

\uline{Démonstration} : Soit $D_s$ un disque de centre $s$ satisfaisant b) de la def. 1 par rapport a $\Lambda^{1}$ \deleted{et}, à $\Lambda^{2}$ et \={a} $\Lambda$.

Observons d'abord que le secteur angulaire exterieur \inserted{\tikz[baseline=(n.base)]{\node[inner sep=2pt] (n) {$D_s\cap\mathrm{ext}\,\Lambda$}; \draw (n.west) .. controls (n.north west) and (n.north east) .. (n.east) .. controls (n.south east) and (n.south west) .. cycle;}} determiné par $\Lambda$ vérifie
\[
D_s\cap\mathrm{ext}\,\Lambda\subset D_s\cap\mathrm{ext}\,\Lambda^{1}\quad(\text{resp. }D_s\cap\mathrm{ext}\,\Lambda^{2}).
\]
En effet tout point $x$ de $D_s\cap\mathrm{ext}\,\Lambda$ peut être relié à un point $y$ de $\mathbb{R}^{2}$ tel que
\[
y\notin(\mathrm{int}\,\Lambda^{1}\cup\mathrm{int}\,\Lambda^{2}\cup\mathrm{int}\,\Lambda)\cup(\Lambda^{1}\cup\Lambda^{2}\cup\Lambda)
\]
\deleted{puisque} (on peut toujours trouver un tel $y$), par un arc de polygone $\Lambda_{xy}$ avec $\Lambda_{xy}\cap\Gamma=\varnothing$ \deleted{et} tel que $\Lambda_{xy}\subset\mathrm{ext}\,\Lambda$. Il est clair que $\Lambda_{xy}\cap\Lambda^{i}=\varnothing$ puisque $\Lambda_{xy}\cap{]s,s'[}=\varnothing$ car $ {]s,s'[}\deleted{\unsure{\not\subset}}\subset\mathrm{int}\,\Lambda^{i}\deleted{\unsure{\ne\varnothing|=\varnothing}}$, comme $y\notin\mathrm{int}\,\Lambda^{i}$, $i=1,2$, cela montre que $x\in\mathrm{ext}\,\Lambda^{i}$, $i=1,2$.

Cela permet de prouver que $(\Lambda^{2}\deleted{\unsure{\cap}}\mathbin{\dot{-}}{]s,s'[})\subset\overline{\mathrm{ext}\,\Lambda^{1}}$, mais le fait que $\Lambda^{2}-{]s,s'[}\subset\mathbb{R}^{2}-\Lambda^{1}$ entraîne que $\Lambda^{2}-{]s,s'[}\subset\mathrm{ext}\,\Lambda^{1}$. On montre aussi \inserted{avec le même argument} que $\Lambda^{1}-{]s,s'[}\subset\mathrm{ext}\,\Lambda^{1}$, c.q.f.d.

\uline{Lemme} $\underline{\overline{\mathrm{XIX}}}$. $\mathrm{int}\,\Lambda^{1}\subset\mathrm{ext}\,\Lambda^{2}$, $\mathrm{ext}\,\Lambda^{1}\supset\mathrm{int}\,\Lambda^{2}$.

\uline{Demonstration} : on montre que tout point de $\mathrm{int}\,\Lambda^{1}$ peut être relié\deleted{\unsure{r|e}} à un point de $\mathrm{ext}\,\Lambda^{2}$ par un arc de polygone qui ne coupe pas $\Lambda^{2}$. Comme $\Lambda^{\unsure{4|1|+}}-\{s,s'\}\subset\mathrm{ext}\,\Lambda^{2}$ soit $x\in\Lambda^{1}-\{s,s'\}$ alors \inserted{il \unsure{exìste|existe}} $D_x\subset\mathrm{ext}\,\Lambda^{2}$ satisfaisant b) de la def. 1. par rapport a $\Lambda^{1}$, donc $D_x\cap\mathrm{int}\,\Lambda^{1}\ne\varnothing$. Soi\inserted{en}t $z\in D_x\cap\mathrm{int}\,\Lambda^{1}$ et $y\in\mathrm{int}\,\Lambda^{1}$ il \unsure{exìste|existe} un arc de polygone $\Lambda_{zy}\subset\mathrm{int}\,\Lambda^{1}$ réliant $z\,(\in\mathrm{ext}\,\Lambda^{2})$ à $y$.
% archive sheet number: 44

\scanpage{46}
\figure{\begin{tikzpicture}[x=.021cm,y=-.021cm,line join=round,line cap=round]
% Upper drawing. Its lower vertex and the second drawing share this coordinate system.
\draw[black!78,line width=.85pt]
  (771,1187) .. controls (815,1219) and (864,1258) .. (908,1324)
  -- (1028,1205) -- (969,1235) -- (1029,1127) -- (1130,1253)
  -- (886,1496) .. controls (857,1397) and (809,1279) .. (771,1187);
\draw[black!63,line width=.55pt]
  (992,1240)--(1025,1207)--(978,1231);
\draw[black!46,line width=.55pt]
  (777,1187)--(795,1188) (812,1190)--(830,1190)
  (875,1195)--(892,1195) (936,1200)--(955,1201)
  (973,1203)--(991,1204) (1001,1204)--(1014,1205);
\draw[black!20,line width=.28pt]
  (908,1326) .. controls (906,1384) and (897,1455) .. (886,1496);
\draw[black!70,line width=1.1pt]
  (908,1328)--(909,1345) (907,1358)--(905,1377)
  (903,1392)--(903,1408) (899,1422)--(896,1443)
  (892,1464)--(887,1494);
\draw[black!48,thin] (886,1494)--(884,1512);
\node[anchor=south east,inner sep=1pt] at (765,1184) {$s'$};
\node[anchor=north west,inner sep=1pt] at (913,1324) {$s$};
\node[anchor=west,inner sep=0pt] at (1025,1210) {$s''$};
% Lower drawing.
\draw[black!60,line width=.8pt]
  (1028,1536) .. controls (1030,1602) and (1042,1690) .. (1055,1754)
  -- (1205,1535) -- (1076,1624) -- cycle;
\draw[line width=.7pt] (1028,1536)--(1076,1624)--(1089,1621);
\draw[black!45,line width=.42pt] (1080,1618)--(1205,1535);
\draw[black!27,line width=.35pt] (1030,1538)--(1038,1604);
\draw[black!48,line width=.42pt] (1028,1536)--(1036,1538);
% Isolated short ink stroke well below the four-cornered outline.
\draw[black!58,line width=.48pt]
  (1114,2020) .. controls (1115,2023) and (1113,2027) .. (1114,2032);
\end{tikzpicture}}

% The writing is inverted (180 degrees) at the lower right.
\margin{ment\qquad $\deleted{\unsure{]}}[s,s']\cap$}
% archive sheet number: 45

\scanpage{47}\mspage{$\overline{\underline{\mathrm{XVIII}}}$}
mais $\Lambda_{zy}\cap{]s,s'[}=\varnothing$ puisque $\Lambda_{zy}\subset\mathrm{int}\,\Lambda^1$, et par la même raison $\Lambda_{zy}\cap(\Lambda^2-\{s,s'\})=\varnothing$ car $\Lambda^2-\{s,s'\}\subset\mathrm{ext}\,\Lambda^1$ (cf. lemme $\overline{\underline{\mathrm{XVIII}}}$) cela prouve la premiēre observation, on obtient la seconde en changeant $\Lambda^1$ par $\Lambda^2$ dans la preuve ci-dessus, c.q.f.d.

\uline{Lemme} $\overline{\underline{\mathrm{XX}}}$. Il y'a deux bons sommets de $\Gamma$, \inserted{\tikz[baseline=(ann.base)]{
\node[inner sep=2pt] (ann) {non consécutifs si $p=$ nb de \unsure{cotés|côtés} de $\Gamma\geq4$,};
\draw ([xshift=-2pt]ann.north west) .. controls ([xshift=-6pt]ann.west) and ([xshift=-5pt]ann.south west) .. (ann.south west) -- (ann.south east) .. controls ([xshift=5pt]ann.south east) and ([xshift=5pt]ann.east) .. ([xshift=2pt]ann.north east);
\draw ([xshift=18pt]ann.south west) -- ++(0,-4pt);
}} \deleted{au moins} tels que les \inserted{\tikz[baseline=(ins.base)]{\node[inner sep=1pt] (ins) {interieurs};\draw (ins.south west) -- ([xshift=-2pt]ins.south) -- ([yshift=-3pt]ins.south) -- ([xshift=2pt]ins.south) -- (ins.south east);} des} triangles corespondants soient contenus dans $\mathrm{int}\,\Gamma$.

\uline{Démonstration} : L'hypothēse de récurrence appliquée à $\Lambda^1$ et $\Lambda^2$ montre qu'il'y a un \uline{bon sommet}\margin{\unsure{tel que}\ldots} % slanted
$t_1$ dans $\Lambda^1$ et un bon sommet $t_2$ dans $\Lambda^2$. Si $t_1\ne s,s'$ alors $t_1$ est un bon sommet de $\Gamma$ puisque $\mathrm{int}\,\Lambda^1\subset\mathrm{int}\,\Lambda$. (tout point de $\mathrm{int}\,\Lambda^1$ peut être relié à un point de $\mathrm{int}\,\Lambda$, par example en prenant un disque $D_x$, $x\in\Gamma^{1}_{ss'}-\{s,s'\}$ vérifiant b) de la def.1 par rapport à $\Gamma$ et à $\Lambda^1$ on voit que $\mathrm{int}\,\Lambda^1\cap\mathrm{int}\,\Lambda\ne\varnothing$....etc)

Dans le cas contraire sans perte de généralité on peut poser $t_1=s$. Deux cas \margin{?}peuvent se produire $ {]\text{\textcircled{\scriptsize \(s''\)}},s'[}\cap\Gamma=\varnothing$ ou bien $ {]s'',s'[}\cap\Gamma\ne\varnothing$ et dans ce cas $ {]s'',s'[}\cap S(\Gamma)\ne\varnothing$ puisque $ {]s'',s'[}\not\subset\Gamma$.

\margin{\deleted{\unsure{faux}}} % slanted
Si $\deleted{\unsure{[}}{]s'',s'[}\cap\Gamma=\varnothing$ il est clair que $s''$ est visible de $s'$ par l'interieur (car le triangle de sommets $s,s',s''\subset\mathrm{int}\,\Gamma$). Considerons alors $\unsure{\Lambda^{1'}_{\mkern5mu\cdot}|\Lambda^{1!}|\Lambda^{11}_{\mkern5mu\cdot}}$ le polygone \margin{!}obtenu en récollant selon \deleted{\illegible{1}} $\{s',s''\}$ l'arc de polygone $[s',s'']$ avec le seul des arcs de polygone de $\Lambda^1$ réliant $s$ à $s'$ qui ne contient pas $[s,s']$. Il est clair que $S(\Lambda^{1'})\subset S(\Gamma)$ et que l'ensemble de toutes les arêtes de $\Lambda^{1'}$ sauf une, savoir $[s',s'']$ est contenu dans $A(\Gamma)$, \unsure{Mais|mais} nécessairement \margin{?}l'ensemble des arètes de $\Lambda^{1'}$ contenues dans $A(\Gamma)$ est de cardinal inferieur \margin{so what} % slanted
au cardinal de l'ensemble des arêtes de $\Lambda^1$ contenues dans $A(\Gamma)$.

Si $ {]s'',s'[}\cap S(\Gamma)\ne\varnothing$ on choisit $z\in S(\Gamma)\cap{]s'',s'[}$, il est clair que $z$ est visible \inserted{de $s=t$} par l'interieur (puisque $ {]s,z[}\subset\overset{\circ}{\text{triangle}}$ de sommets $s,s',s''$). à partir du couple $s,z$ on peut appliquer la même construction que ci-dessus pour aboutir aussi à un polygone $\Lambda^{1'}$ tel que : a) $S(\Lambda^{1'})\subset S(\Gamma)$; b)
% archive sheet number: 46

\scanpage{48}\mspage{$\overline{\underline{\mathrm{XIX}}}$}
tou\deleted{s}tes les \deleted{\unsure{so|sa}}\unsure{.} arêtes de $\Lambda^{1'}$ sauf une (savoir, $[z,s]$) appartiennent à $\deleted{\unsure{R|A}}\,A(\Gamma)$ et le cardinal de vérifie $A(\Lambda^{1'})\cap A(\Gamma)<A(\Gamma)\cap A(\Lambda^1)$.

Comme il n'ya qu'un nombre fini d'arêtes de $\Gamma$ il est clair qu'on \deleted{\unsure{bien}} a :

a) Il existe $n$ tel que $\deleted{\unsure{A|R}}\,\Lambda^{1(n)}$ (Notation definie par récurrence de foçon evidente) contienne un \inserted{bon} sommet $t$ tel que les deux arêtes qui contiennent $t$ soient des arêtes $\Gamma$, et par conséquent $t$ est aussi un bon sommet de $\Gamma$ compte tenu du lemme $\overline{\underline{\mathrm{XIX}}}$, appliqué $n$-fois, \inserted{ou bien,}

b) il existe $n$ tel que $\mathrm{card}\,(A({\Lambda}^{1(n)})\cap A(\Gamma))=2$ donc $\mathrm{card}\,A^{1(n)}=3$ et il est clair que le \inserted{seul} sommet de $\Lambda^{1(n)}$ qui appartient aux deux arêtes de $\Lambda^{1(n)}$ contenues dans $A(\Gamma)$ est un bon sommet de $\Gamma$, compte tenu encore une fois du lemme $\overline{\underline{\mathrm{XIX}}}$\inserted{\tikz[baseline=(a.base)]{
\node[inner sep=1pt] (a) {appliqué $n$ fois Le même argument avec $\Lambda^2$};
\node[anchor=west,inner sep=1pt] (b) at ([xshift=120pt,yshift=-11pt]a.south west) {prouve le reste,};
\draw ([xshift=-4pt,yshift=-4pt]a.south west) -- (a.south west) -- ([xshift=115pt]a.south west) .. controls ([xshift=130pt,yshift=1pt]a.south west) and ([xshift=-8pt,yshift=7pt]b.south west) .. (b.south west) -- ([xshift=10pt]b.south east);
}}, c.q.f.d.

\uline{Lemme $\overline{\underline{\mathrm{XXI}}}$} : Si $\Gamma$ est un polygone non-convexe il existe un bon sommet de $\Gamma$ tel que le \inserted{interieur du} triangle corespondant est contenu dans $\mathrm{ext}\,\Gamma$.

\uline{Demonstration} : le lemme $\overline{\underline{\mathrm{VIII}}}$ montre qu il \unsure{exìste|existe} un sommet $s$ de $\Gamma$ et un sommet $s'$ visible de $s$ par l'exterieur de $\Gamma$. On conserve ici les notations du \unsure{numèro|numéro|numero} $3$. On considère $\Lambda^2$ (defini comme dans le numéro $3$) il est clair \margin{\tikz[baseline=(m.base)]{\node[align=left,inner sep=1pt] (m) {faux, si on\\ne choisit\\pas bien\\\unsure{s.|ses} \unsure{notations|solutions} \unsure{--|…}};\draw ([xshift=2pt]m.north east) -- ([xshift=-5pt]m.south east);}} % slanted
que $\mathrm{int}\,\Lambda^2\subset\mathrm{ext}\,\Gamma$ car $V=\mathrm{ext}\,\Lambda^1\cup\mathrm{int}\,\Lambda^2\cup{]s,s'[}$, et $V$ est non \deleted{\unsure{co|lo}}-relativement compact. \deleted{L'argument du}\inserted{Le} lemme précedent appliqué a $\Lambda^2$ montre que l'on peut trouver un \deleted{\unsure{trian|triang}} bon sommet $s$ de $\Lambda^2$ tel que : le triangle corespondent soit contenu dans $\mathrm{ext}\,\Gamma$ et les deux arêtes de $\Lambda^2$ auxquelles \margin{\uline{\unsure{douteux|donc ?}}} % slanted
$s$ appartient soient des arêtes de $\Gamma$, ce qui permet de prouver le lemme.

\deleted{\unsure{Le|Les}}\unsure{Le|Lee} lemme $\overline{\underline{\mathrm{XX}}}$ le lemme $\overline{\underline{\mathrm{XXI}}}$ joints au fait que dans tout polygone convexe il 'y a au moins trois bons sommets, \unsure{montent|montrent} le théorème $2$, c.q.f.d.

\begin{flushright}
le 15 juillet 1976, Montpellier\\
\unsure{Mauré|Mouré|Maury|Maurice}.
\end{flushright}
% archive sheet number: 47

